% =====================================================================
%  FICHE 4 — THEME 1 — Exercices niveau FACILE
% =====================================================================
\documentclass[11pt,a4paper]{article}
\usepackage[utf8]{inputenc}
\usepackage[T1]{fontenc}
\usepackage{lmodern}
\usepackage[french]{babel}
\usepackage{amsmath,amssymb,mathtools}
\usepackage{icomma}
\usepackage[version=4]{mhchem}
\usepackage{siunitx}
\sisetup{output-decimal-marker={,},per-mode=symbol}
\DeclareSIUnit\Molar{M}
\usepackage{xcolor}
\usepackage{tikz}
\usepackage[most]{tcolorbox}
\usepackage{enumitem}
\usepackage{booktabs}
\usepackage{array}
\usepackage{fancyhdr}
\usepackage[a4paper,margin=2.2cm,headheight=26pt,headsep=10pt]{geometry}

\definecolor{primary}{RGB}{146,95,160}
\definecolor{primarydark}{RGB}{92,52,104}
\definecolor{defcol}{RGB}{0,114,178}
\definecolor{formulacol}{RGB}{198,93,9}

\tcbset{boxbase/.style={enhanced,breakable,boxrule=0.9pt,arc=2.5pt,
   left=3mm,right=3mm,top=2mm,bottom=2mm,fonttitle=\bfseries,coltitle=white,
   toptitle=0.5mm,bottomtitle=0.5mm}}
\newtcolorbox{donnees}{boxbase,colback=defcol!5,colframe=defcol,
   title={\textsc{Données}}}

\newcounter{exo}
\newenvironment{exo}[1][]{%
  \par\medskip\refstepcounter{exo}%
  \noindent\begin{tcolorbox}[boxbase,colback=primary!7,colframe=primary,
     title={\textsc{Exercice \theexo}\ifx\relax#1\relax\relax\else\textmd{ --- \itshape #1}\fi}]}%
  {\end{tcolorbox}}

\pagestyle{fancy}\fancyhf{}
\fancyhead[L]{\small\textcolor{primarydark}{\textbf{Fiche 4} --- Th.\,1 : Mole, masse molaire, entit\'es}}
\fancyhead[R]{\small\textcolor{primarydark}{\textsc{Chimie} --- \textbf{Facile}}}
\fancyfoot[C]{\small\thepage}
\renewcommand{\headrulewidth}{0.8pt}
\renewcommand{\headrule}{\hbox to\headwidth{\color{primary}\leaders\hrule height \headrulewidth\hfill}}
\newcommand{\NA}{N_{\!\text{A}}}
\setlength{\parindent}{0pt}

\begin{document}

\begin{center}
{\LARGE\bfseries\color{primarydark} Thème 1 --- Niveau Facile}\\[2pt]
{\color{primary}\itshape Exercices : compter et peser la matière}
\end{center}

\begin{donnees}
Masses molaires atomiques (\si{\gram\per\mole}) :
$\ce{H}=1$ ; $\ce{C}=12$ ; $\ce{N}=14$ ; $\ce{O}=16$ ; $\ce{Na}=23$ ; $\ce{S}=32$ ;
$\ce{Fe}=56$ ; $\ce{Cu}=63{,}5$.\quad
Nombre d'Avogadro : $\NA=\SI{6.02e23}{\per\mole}$.
\end{donnees}

\begin{exo}[masse molaire d'un composé]
Calculer la masse molaire des composés suivants, en détaillant la somme atome par atome :
\begin{enumerate}[label=\textbf{\alph*)},leftmargin=2em,topsep=2pt,itemsep=2pt]
  \item le dioxyde de carbone \ce{CO2} ;
  \item le glucose \ce{C6H12O6} ;
  \item l'urée \ce{CH4N2O} ;
  \item le sulfate de fer(II) \ce{FeSO4}.
\end{enumerate}
\end{exo}

\begin{exo}[de la masse à la quantité de matière]
On pèse $\SI{20}{\gram}$ d'hydroxyde de sodium \ce{NaOH}, dont la masse molaire vaut
$M(\ce{NaOH})=\SI{40}{\gram\per\mole}$.\\
Quelle \textbf{quantité de matière} (en moles) de \ce{NaOH} cet échantillon représente-t-il ?
\end{exo}

\begin{exo}[de la quantité de matière au nombre d'entités]
On dispose de $\SI{0.25}{\mole}$ de molécules d'eau \ce{H2O}.
\begin{enumerate}[label=\textbf{\alph*)},leftmargin=2em,topsep=2pt,itemsep=2pt]
  \item Combien de \textbf{molécules} d'eau cela représente-t-il ?
  \item Combien de \textbf{moles d'atomes} d'hydrogène y a-t-il dans cet échantillon ?
\end{enumerate}
\end{exo}

\begin{exo}[du nombre d'entités à la quantité de matière]
Un flacon contient $\num{3.01e23}$ molécules de dioxyde de carbone \ce{CO2}.
\begin{enumerate}[label=\textbf{\alph*)},leftmargin=2em,topsep=2pt,itemsep=2pt]
  \item Quelle quantité de matière de \ce{CO2} (en moles) contient le flacon ?
  \item En déduire la quantité de matière d'\textbf{atomes d'oxygène} correspondante.
\end{enumerate}
\end{exo}

\begin{exo}[masse molaire d'un hydrate]
Certains sels cristallisent avec des molécules d'eau, dites « eau de cristallisation ». Calculer la
masse molaire des sels hydratés suivants (on rappelle $M(\ce{H2O})=\SI{18}{\gram\per\mole}$) :
\begin{enumerate}[label=\textbf{\alph*)},leftmargin=2em,topsep=2pt,itemsep=2pt]
  \item le sulfate de cuivre(II) pentahydraté \ce{CuSO4.5H2O} ;
  \item le carbonate de sodium décahydraté \ce{Na2CO3.10H2O}.
\end{enumerate}
\end{exo}

\end{document}
